\section{Algebraic Preliminaries}
\label{sec:algebra-preliminaries}

\subsection{Finite fields}

Most SNARKs arithmetize computation over a finite field. In these notes we usually write
\[
    \F=\F_p
\]
for a prime field. Elements are integers modulo a prime $p$, and every nonzero element has a multiplicative inverse. The field structure lets arithmetic circuits use addition and multiplication as primitive operations, while the finiteness of the field makes random evaluation arguments meaningful.

The field must be large compared with the degrees of the polynomials that appear in the proof. A typical soundness statement has the form
\[
    \Pr[\text{bad polynomial identity passes}]\le \frac{D}{p},
\]
where $D$ is a degree bound. Cryptographic systems therefore choose $p$ large enough that this ratio is negligible.

\subsection{Univariate polynomials, roots, and degree}

A univariate polynomial over $\F$ is
\[
    f(X)=f_0+f_1X+\cdots+f_dX^d.
\]
The largest $i$ with $f_i\neq 0$ is the degree. The elementary fact used throughout these notes is:
\[
    \text{a nonzero polynomial of degree }D\text{ has at most }D\text{ roots.}
\]
This is the algebraic reason that checking an identity at a random point can certify a global identity with high probability.

\subsection{Evaluation and interpolation}

A degree-$<k$ polynomial is uniquely determined by its values at $k$ distinct field points. Given distinct points $a_0,\ldots,a_{k-1}$ and values $y_0,\ldots,y_{k-1}$, the interpolating polynomial is
\[
    f(X)=\sum_{i=0}^{k-1} y_i\lambda_i(X),
\]
where the Lagrange basis polynomial is
\[
    \lambda_i(X)=\prod_{j\neq i}\frac{X-a_j}{a_i-a_j}.
\]
Thus SNARKs can freely switch between coefficient representations and evaluation representations. Plonkish systems strongly prefer the evaluation viewpoint: witness values are placed on a structured domain, then interpolated into low-degree polynomials.

\subsection{Roots of unity and subgroup domains}

Let $\omega\in\F$ be a primitive $k$-th root of unity. Then
\[
    \omega^k=1,
    \qquad
    \omega^i\neq 1\quad\text{for }0<i<k.
\]
The associated multiplicative subgroup is
\[
    \Omega=\{1,\omega,\omega^2,\ldots,\omega^{k-1}\}.
\]
Such domains are central because shifting by $\omega$ moves to the next row:
\[
    X\mapsto \omega X.
\]
This is exactly what ProductCheck and Plonk gate checks need: local relations between neighboring positions become polynomial identities involving $f(X)$, $f(\omega X)$, and sometimes $f(\omega^2X)$.

\subsection{Vanishing polynomials}

For a finite set $\Omega\subseteq\F$, the vanishing polynomial is
\[
    Z_\Omega(X)=\prod_{a\in\Omega}(X-a).
\]
A polynomial $f$ vanishes on $\Omega$ iff $Z_\Omega$ divides $f$:
\[
    f(a)=0\ \forall a\in\Omega
    \quad\Longleftrightarrow\quad
    f(X)=q(X)Z_\Omega(X).
\]
For a multiplicative subgroup $\Omega=\{1,\omega,\ldots,\omega^{k-1}\}$,
\[
    Z_\Omega(X)=X^k-1.
\]
This special form is why subgroup domains are so efficient: the verifier can evaluate $Z_\Omega(r)=r^k-1$ in $O(\log k)$ field operations.

For arbitrary subsets, the same theory works with
\[
    Z_\Omega(X)=\prod_{a\in\Omega}(X-a),
\]
but the verifier no longer gets the simple $X^k-1$ formula unless the subset is itself a subgroup or has some other special structure.

\subsection{Polynomial identity testing}

The DeMillo--Lipton--Schwartz--Zippel principle says that a nonzero polynomial of total degree $D$ over a field cannot vanish too often \cite{DeMilloLipton1978,Schwartz1980,Zippel1979}. In the univariate case,
\[
    \Pr_{r\leftarrow \F}[f(r)=0]\le \frac{\deg f}{|\F|}
\]
for nonzero $f$. More generally, for a nonzero multivariate polynomial $F$ of total degree $D$,
\[
    \Pr_{\mathbf r\leftarrow S^m}[F(\mathbf r)=0]\le \frac{D}{|S|}.
\]
This principle is the soundness engine behind polynomial equality testing, ZeroTest, ProductCheck, and permutation arguments.

\subsection{Cyclic groups, discrete logarithms, and pairings}

Polynomial commitments are cryptographic objects, so they require group assumptions in addition to field algebra. Let $\G=\langle g\rangle$ be a cyclic group of prime order $p$. In multiplicative notation, every group element has the form $g^a$ for some $a\in\F_p$, and
\[
    g^a\cdot g^b = g^{a+b},
    \qquad
    (g^a)^b=g^{ab}.
\]
The discrete logarithm problem is: given $g$ and $h=g^x$, find $x$. Pedersen commitments and Bulletproofs-style commitments rely on the hardness of recovering or manipulating hidden linear relations among such exponents \cite{Pedersen1991,Bulletproofs2018}.

Many pairing-based SNARKs also use a bilinear map
\[
    e:\G\times\G\to\GT
\]
with
\[
    e(g^a,g^b)=e(g,g)^{ab}.
\]
Pairings let a verifier test multiplicative relations among hidden exponents. A classical cryptographic example is the BLS short-signature construction~\cite{BonehLynnShacham2001}. KZG and Groth16 both rely on this ability, though in different ways \cite{KateZaveruchaGoldberg2010,Groth2016}.

